% Part: turing-machines
% Chapter: machines-computations
% Section: halting-states

\documentclass[../../../include/open-logic-section]{subfiles}

\begin{document}

\olfileid[pt]{tur}{mac}{hal}
\olsection{Estados de Parada}

\begin{explain}
Embora tenhamos definido nossas máquinas para parar apenas quando
não há instrução a executar, representações comuns de máquinas de
Turing têm um \emph{estado de parada} dedicado~$h$, de modo que
$h \in Q$.

A ideia por trás de um estado de parada é simples: quando a máquina
terminou a operação (está pronta para aceitar a entrada, ou terminou
de escrever a saída), ela entra em um estado~$h$ onde para. Algumas
máquinas têm dois estados de parada, um que aceita a entrada e outro que
rejeita a entrada.
\end{explain}

\begin{ex}\emph{Estados de parada}.
Para elucidar esse conceito, comecemos com uma alteração da
máquina dos pares. Em vez de fazer a máquina parar no estado~$q_0$
se a entrada for par, podemos adicionar uma instrução para enviar a máquina
a um estado de parada.
\[
\begin{tikzpicture}[->,>=stealth',shorten >=1pt,auto,node distance=2.8cm,
                    semithick]
  \tikzstyle{every state}=[fill=none,draw=black,text=black]

  \node[initial, state]         (A)                     {$q_0$};
  \node[state]         (B) [right of=A] {$q_1$};
  \node[state]         (C) [below of=A] {$h$};

  \path (A) edge [bend left] node {\TMtrans{\TMstroke}{\TMstroke}{R}} (B)
  	    edge node {\TMtrans{\TMblank}{\TMblank}{N}} (C)
        (B) edge [loop above] node {\TMtrans{\TMblank}{\TMblank}{R}} (B)
            edge [bend left] node {\TMtrans{\TMstroke}{\TMstroke}{R}} (A);
\end{tikzpicture}
\]

Vamos expandir ainda mais o exemplo. Quando a máquina determina que a entrada
é ímpar, ela nunca para. Podemos alterar a máquina para
incluir um estado de \emph{rejeição} substituindo a instrução em laço por uma
instrução para ir a um estado de rejeição~$r$.
\[
\begin{tikzpicture}[->,>=stealth',shorten >=1pt,auto,node distance=2.8cm,
                    semithick]
  \tikzstyle{every state}=[fill=none,draw=black,text=black]

  \node[initial,state]         (A)                     {$q_0$};
  \node[state]         (B) [right of=A] {$q_1$};
  \node[state]         (C) [below of=A] {$h$};
  \node[state]         (D) [below of=B] {$r$};

  \path (A) edge [bend left] node {\TMtrans{\TMstroke}{\TMstroke}{R}} (B)
  	    edge node {\TMtrans{\TMblank}{\TMblank}{N}} (C)
        (B) edge node {\TMtrans{\TMblank}{\TMblank}{N}} (D)
            edge [bend left] node {\TMtrans{\TMstroke}{\TMstroke}{R}} (A);
\end{tikzpicture}
\]
\end{ex}

\begin{explain}
Adicionar um estado de parada dedicado pode ser vantajoso em casos como
este, em que deixa explícito quando a máquina aceita/rejeita certas
entradas. No entanto, é importante notar que nenhum poder computacional é
ganho ao adicionar um estado de parada dedicado. Da mesma forma, uma noção
menos formal de parada tem suas próprias vantagens. A definição de
parada usada até agora neste capítulo torna a prova do
\emph{Problema da Parada} intuitiva e fácil de demonstrar. Por essa
razão, continuamos com nossa definição original.
\end{explain}

\end{document}
